\begin{claim}{D.29}{Refuted}{A preserved four-puncture structure on the torsion-free Haran extension}
The full profinite recognition assertion in D(d) is false, including
its peripheral hypothesis.  There is an exact sequence of profinite
groups
\[
 1\longrightarrow H=\widehat\pi_{0,4}
 \longrightarrow E\longrightarrow C_2\longrightarrow1
                                                        \tag{D.29.1}
\]
in which $E$ is torsion-free and its faithful outer action preserves
a complete type-$(0,4)$ loop structure, while $E$ is the completion
of no finite-type surface group, orientable or nonorientable.
The extension is explicitly isomorphic to D.26's group.
Its outer involution has a mixed-prime orientation multiplier and
is conjugate to no subgroup of the geometric completion inside
$\Out_{\rm lac}(\widehat\pi_{0,4})$.
For this involution both the discretization and prediscretization
fixed sets of F.1 are empty.
\end{claim}
\begin{proof}
Define units and an exponent by their prime coordinates:
\[
 \epsilon_3=-1,\qquad\epsilon_q=1\ (q\ne3),\qquad
 \kappa_2=1,\quad\kappa_q=0\ (q\ne2),\qquad
 \lambda=-\epsilon.
                                                        \tag{D.29.2}
\]
Thus $\epsilon^2=\lambda^2=1$ and
$\epsilon\kappa=\kappa$.  In particular $\lambda_3=1$
and $\lambda_q=-1$ for every other prime.

We first obtain the required pants isomorphism, without any
finite-order requirement on it.  The compatible assignments
$\zeta_n\mapsto\zeta_n^{\lambda\bmod n}$ define an automorphism
of the cyclotomic extension $\Q^{\rm cyc}/\Q$.
An automorphism of an algebraic extension extends to its algebraic
closure: extend its embedding along algebraic generators, or take
a maximal extension of the embedding and adjoin a root of a missing
minimal polynomial to prove maximality gives all of
$\overline\Q$.  We therefore choose an element
$\sigma\in G_{\Q}$ with cyclotomic character $\lambda$.
The Galois action on the three-punctured projective line commutes
outerly with its rational $S_3$ action and preserves each cusp;
on its local inertia it has exponent $\lambda$.
Harbater--Schneps's normalized description
\cite[\S0.2, Main Theorem (a), Theorem A]{R-HS}
therefore supplies an element $(\lambda,f)\in GT_0$, with
$f\in\widehat F(x,y)'$, whose automorphism is
\[
 x\longmapsto x^\lambda,\qquad
 y\longmapsto f(x,y)^{-1}y^\lambda f(x,y).
                                                        \tag{D.29.3}
\]
Its third cusp $(xy)^{-1}$ goes to a conjugate of its
$\lambda$th power.  This is the only property of $f$ needed
below, besides invertibility of the displayed automorphism.
Every choice of the extension $\sigma$ has these properties,
and every one gives the construction and conclusions that follow.

Let $H=\widehat F(l,a,b)$ and define its closed free factors
\[
 P_1=\widehat F(l,a),\qquad P_2=\widehat F(l,b),\qquad
 L=\overline{\langle l\rangle}.
\]
Their shared-factor universal property is
$H=P_1\mathbin{*_{L}}P_2$: homomorphisms out of the two
free factors agreeing on $l$ are exactly arbitrary continuous
assignments of $l,a,b$.  Give $P_1$ the ordered pants basis
$(l^{-1},a)$ and $P_2$ the ordered pants basis $(l,b)$.
Formula (D.29.3), between these bases, defines a continuous
isomorphism $\tau:P_1\to P_2$ with
\[
 \begin{split}
 \tau(l)&=l^\epsilon,\\
 \tau(a)&=f(l,b)^{-1}b^\lambda f(l,b),\\
 \tau(a^{-1}l)&\sim_{P_2}(b^{-1}l^{-1})^\lambda.
 \end{split}                                           \tag{D.29.4}
\]
The symbol $\sim_{P_2}$ means conjugacy by an element of
$P_2$; in particular it specifies a conjugacy in $H$.
Because $\lambda^{-1}=\lambda$, the inverse isomorphism has
the same peripheral multiplier.

Put $c=l^\kappa$ and prescribe
\[
 \alpha|_{P_1}=\tau,\qquad
 \alpha|_{P_2}=\operatorname{Inn}(c)\tau^{-1}.
                                                        \tag{D.29.5}
\]
Both maps send the shared generator $l$ to $l^\epsilon$:
$\tau^{-1}(l)=l^\epsilon$ and $c$ commutes with $l$.
They consequently extend to a continuous endomorphism of
$H$.  Direct substitution on the two free factors gives
\[
                  \alpha(c)=c,\qquad
                  \alpha^2=\operatorname{Inn}(c).      \tag{D.29.6}
\]
For example, on $P_1$ the second iterate is
$\operatorname{Inn}(c)\tau^{-1}\tau$; on $P_2$ it is
$\operatorname{Inn}(\alpha(c))\tau\tau^{-1}$.
These identities also prove invertibility, with inverse
$\operatorname{Inn}(c^{-1})\alpha$.

Here is the full peripheral structure, with its relation:
\[
 \begin{split}
 (u_1,u_2,u_3,u_4)&=(a,\ a^{-1}l,\ b,\ b^{-1}l^{-1}),\\
                         u_1u_2u_3u_4&=1.
 \end{split}                                           \tag{D.29.7}
\]
The first three entries form a free profinite basis, since
$a=u_1$, $l=u_1u_2$, and $b=u_3$ are its continuous
inverse substitutions.  Thus (D.29.7) is precisely a
type-$(0,4)$ presentation, not just four procyclic subgroups
whose generators happen to satisfy a relation.
Equations (D.29.4)--(D.29.5) give
\[
 \alpha(u_i)\sim_H u_{s(i)}^\lambda,
                  \qquad s=(1\ 3)(2\ 4).              \tag{D.29.8}
\]
For the first two entries this is (D.29.4), and for the last
two it is the inverse comparison, followed by conjugation
by $c$.  Hence the action preserves every required loop
source, its conjugacy ambiguity, and the common orientation
multiplier.  It permutes distinct cusp classes nontrivially,
so its outer class is nonidentity.  Equation (D.29.6) gives
that class exact order two.

On the profinite space $H\times\{0,1\}$ define
\[
 (h,i)(k,j)=
   \bigl(h\alpha^i(k)c^{ij},\ i+j\bmod2\bigr).
                                                        \tag{D.29.9}
\]
Associativity follows by the two carry identities in
(D.29.6); checking the binary indices gives precisely
$\alpha^2(k)c=ck$ and $\alpha(c)=c$.
The element $t=(1,1)$ satisfies $t^2=c$ and
$tht^{-1}=\alpha(h)$.  The inclusion $h\mapsto(h,0)$
and parity projection give (D.29.1), with the profinite
product topology and continuous operations.

We make the comparison with D.26 explicit.  Set
$b'=\tau(a)$.  The pair $(l^\epsilon,b')$ is a free basis
of $P_2$, and replacing its first member by its
$\epsilon$th power gives the free basis $(l,b')$.
Thus $(l,a,b')$ is a free basis of $H$, and
\[
 \alpha(l)=l^\epsilon,\qquad
 \alpha(a)=b',\qquad
 \alpha(b')=c a c^{-1}.                               \tag{D.29.10}
\]
The map from D.26's named free basis to $(l,a,b')$, with
$t\mapsto t$, is a continuous isomorphism of its entire
twisted extension onto $E$.  It is an isomorphism of
extensions with identity map on $C_2$.  Formula (D.29.7)
therefore supplies the previously missing preserved
peripheral system on that exact group.  In particular
D.26's necessary double-transposition and double-cover
character conditions are attained.

For clarity, we give the torsion and finite-quotient
arguments also in the present basis.  The retraction
$r:H\to\widehat\Z$ sends $l$ to $1$ and $a,b$ to $0$.
Since both $\tau(a)$ and $\tau^{-1}(b)$ lie in the
normal closures of the respective nonshared generators,
$r\alpha=\epsilon r$.  It extends to a retraction of
$E$ onto the group $N=\widehat\Z\times\{0,1\}$ with law
(D.29.9), action $\epsilon$, and square parameter
$\kappa$.  An element $(z,1)$ in $N$ has square
$((1+\epsilon)z+\kappa,0)$, whose $2$-coordinate is
$2z_2+1$ and is never zero.  Since its open subgroup
$\widehat\Z$ is torsion-free, every torsion element
outside it would have order two; hence $N$ is torsion-free.
A torsion element of $E$ has trivial image in $N$, so
has even parity and lies in the torsion-free $H$.
It is therefore the identity.

Let $V=\F_3^2$, with basis $v_1,v_2$, and send
\[
         l\longmapsto v_1,\qquad
         a\longmapsto v_2,\qquad
         b\longmapsto -v_2.
                                                        \tag{D.29.11}
\]
Here $\lambda_3=1$ and $\epsilon_3=-1$; conjugators
vanish in the abelian target.  The induced action of
$\alpha$ is $-1$ on $V$, and $c$ has image zero because
$\kappa_3=0$.  Sending $t$ to the involution consequently
extends (D.29.11) to a continuous epimorphism
\[
              E\longrightarrow V\rtimes_{-1}C_2.      \tag{D.29.12}
\]
The target has order eighteen and requires three
generators.  Two positive elements do not generate the
quotient $C_2$; one positive and one negative element
generate at most one line of $V$; two negative elements
reduce to the preceding case by taking their product.
The three displayed images do generate it.
The independent program
\path{haran_quotient_independent_audit.py} checks all
$18^3$ associativity cases, all $18^2$ ordered pairs,
all ordered triples, and the actual matrix identity in
(D.29.11).  The ordered-pair subgroup-size histogram is
\[
 \begin{array}{c|rrrrr}
 \text{generated subgroup order}&1&2&3&6&9\\\hline
 \text{ordered pairs}&1&27&32&216&48.
 \end{array}                                          \tag{D.29.13}
\]
These calculations agree with the independently executed
D.26 certificate under (D.29.10).

A free profinite $E$ of finite rank $d$ would have an
index-two free subgroup of rank $1+2(d-1)$.  The actual
rank of $H$ is three, forcing $d=2$, contrary to
(D.29.12).  Thus $E$ is not free.  This excludes all
punctured finite-type surface completions.
For a putative isomorphism from a closed surface
completion to $E$, the inverse image of $H$ is an open
index-two subgroup.  Its intersection with the dense
surface group has index two, and completion of this
intersection is exactly that open subgroup by the
finite-index topology argument of D.1.  It is the
completion of the fundamental group of the corresponding
closed double cover.  Every aspherical closed surface
completion has cohomological $p$-dimension two, whereas
$\widehat F_3$ has dimension one.  The spherical and
projective-plane groups and their completions are finite
and cannot give $H$.  This excludes every closed case
as well.

Finally A.2 supplies the continuous orientation
homomorphism
\[
 \operatorname{or}:\Out_{\rm lac}(H)
                       \longrightarrow\widehat\Z^\times.
\]
It is unchanged by conjugation.  Every discrete geometric
class has value $1$ or the global unit $-1$, and the
closure of the geometric group still has values in that
closed two-element subgroup.  The value of $[\alpha]$
is the mixed unit $\lambda$ in (D.29.2), which belongs
to neither value.  No conjugate of $\langle[\alpha]\rangle$
is therefore contained in the geometric completion.
The transporter formulas of D.16 and F.1 now give both
empty fixed sets claimed above.

The cusp normalizers themselves cause no exception to
this argument.  Since $s$ fixes none of the four cusp
classes, their normalizers in $E$ map trivially to $C_2$;
each equals its normalizer in $H$, namely its own
procyclic loop group by A.7.  The failure is present
even with these procyclic normalizers retained.
\end{proof}
